% STATUS: written 2026-11-05 (Chebotar--Ke--Lee 2004, Pan 2012 and
% Brice--Huang 2015 read in full; excerpts in excerpts/. The notion of a zero
% product determined algebra comes from Bresar, Grasic and Sanchez Ortega 2009,
% recorded as atlas-pending; the book names it through Brice--Huang 2015
% pp. 326--327.)
\section{Derivations through zero products and derivable maps}\label{sec:local-characterizations:zero-products}

Sections~\ref{sec:local-characterizations:local}
and~\ref{sec:local-characterizations:two-local} tested the map point by
point and pair by pair. Chebotar, Ke and Lee test it on the zero products
\cite[p.~217]{ChebotarKeLee2004}. An additive map $\delta\colon A\to A$
\emph{behaves like a derivation on the zero products} when
$\delta(x)y+x\delta(y)=0$ holds whenever $xy=0$. Their aim is to get the
derivation back, at least almost \cite[p.~217]{ChebotarKeLee2004}.

The property that governs this section has a name. An algebra $A$ is
\emph{zero product determined} \cite[p.~1487]{BresarGrasicSanchezOrtega2009}
when for every module $X$ and every bilinear map $\{\cdot,\cdot\}$ into $X$,
the condition $\{x,y\}=0$ whenever $xy=0$ forces $\{x,y\}=T(xy)$ for some
linear $T$. Bre\v{s}ar, Gra\v{s}i\v{c} and S\'anchez Ortega introduce it and
prove that the full matrix algebras over a field are zero product determined
\cite[p.~1487]{BresarGrasicSanchezOrtega2009}. A derivable map is the case
where $\{\cdot,\cdot\}$ is built from the map and the product, which is why the
two notions meet here.

The same question is asked of the Jordan product. $A$ is
\emph{zero Jordan product determined} \cite[p.~90]{AnGhahramaniHeidari2015ZeroJordan}
when every bilinear $\varphi$ vanishing on the pairs with $a\circ b=0$ has the
form $\varphi(a,b)=T(a\circ b)$. Write $J(A)$ for the subalgebra
generated by the idempotents. An, Ghahramani and Heidari prove that $A$ is zero
Jordan product determined whenever $A=J(A)$, by exhibiting the map explicitly:
for $a\in A$ and $x\in J(A)$ the value $\varphi(a,x)$ is
$\frac12\varphi(ax,1)+\frac12\varphi(xa,1)$
\cite[Theorem~2.1, p.~91]{AnGhahramaniHeidari2015ZeroJordan}. That settles the
question left open at the end of the paper cited above.

The condition can be weakened to the idempotents alone. $A$ is
\emph{idempotent elements determined} \cite[Definition~1.2, p.~2890]{WangLiChen2011Idempotent}
when the only pairs tested are those with $x$ idempotent, that is when
$\{x,x\}=\{e,x\}$ for every idempotent $x$ already forces $\{x,y\}=f(x\circ y)$.
Wang, Li and Chen prove that $M_{n}(R)$ is idempotent elements determined for a
commutative ring with $2$ invertible \cite[Theorem~2.4]{WangLiChen2011Idempotent}.
So on matrix algebras the whole of the Jordan product need not be tested; the
idempotents carry it.

\begin{theorem}[{\cite[Theorem~2, p.~218]{ChebotarKeLee2004}}]\label{thm:local:cklzero}
Let $A$ be a prime ring with center $Z$, maximal right quotient ring $Q$
and extended centroid $C$. Let $\delta\colon A\to A$ be an additive map such
that $\delta(x)y+x\delta(y)=0$ for all $x,y\in A$ with $xy=0$. Suppose that
$Q$ contains a nontrivial idempotent $e$ such that $eA\cup Ae\subseteq A$.
\begin{enumerate}
\item If $1\in A$, then $\delta(xy)=\delta(x)y+x\delta(y)-\lambda xy$ for
all $x,y\in A$, where $\lambda=\delta(1)\in Z$. In particular, if
$\delta(1)=0$, then $\delta$ is a derivation on $A$.
\item If $\deg A\geq 3$, there exists $\lambda\in C$ such that
$\delta(xy)=\delta(x)y+x\delta(y)-\lambda xy$ for all $x,y\in A$.
\end{enumerate}
\end{theorem}

Chebotar, Ke and Lee also treat the multiplicative side
\cite[p.~217]{ChebotarKeLee2004}. A bijective
additive map of prime rings that preserves zero products is a scalar
multiple of a ring isomorphism
\cite[Theorem~1, pp.~217--218]{ChebotarKeLee2004}.

Pan asks the question at one element
\cite[p.~4251]{Pan2012Derivable}. A linear map $\delta$ of an algebra
$A$ into an $A$-bimodule $M$ is \emph{derivable at $C\in A$} when
$\delta(A)B+A\delta(B)=\delta(C)$ for all $A,B$ with $AB=C$
\cite[p.~4251]{Pan2012Derivable}.

\begin{definition}[{\cite[p.~4251]{Pan2012Derivable}}]\label{def:local:derivpoint}
Let $A$ be a unital algebra and $M$ a unital $A$-bimodule. An element $C\in
A$ is a \emph{derivational point} of $L(A,M)$ when every $\delta\in L(A,M)$
that is derivable at $C$ is a derivation.
\end{definition}

Zero products are the point $C=0$, the Jordan condition is the point $C=I$,
and the recent work asks for the points in between
\cite[p.~4251]{Pan2012Derivable}. Pan describes the maps derivable at a
fixed $C$ through a Peirce decomposition
\cite[Theorem~2.1, p.~4253]{Pan2012Derivable} and then decides which
elements are derivational points.

\begin{theorem}[{\cite[Theorem~4.6, p.~4259]{Pan2012Derivable}}]\label{thm:local:derivpoint}
If $A=\operatorname{alg}\mathcal{N}$ is a nest algebra on a Hilbert space
$H$ then every $0\neq C\in A$ is a derivational point of $L(A,B(H))$.
\end{theorem}

For $C=0$ the condition is exactly the hypothesis of
Theorem~\ref{thm:local:cklzero}; Theorem~\ref{thm:local:derivpoint} says a
nest algebra needs no idempotent hypothesis, since each nonzero element
decides the question by itself. Brice and Huang turn the picture around and
ask which algebras support it \cite[p.~326]{BriceHuang2015}.

\begin{definition}[{\cite[pp.~326--327]{BriceHuang2015}}]\label{def:local:zpd}
Let $K$ be a commutative ring with identity. A $K$-algebra $A$, not
necessarily associative, is \emph{zero product determined} when each
$K$-bilinear map $\varphi\colon A\times A\to B$ with $\varphi(a_1,a_2)=0$
whenever $a_1a_2=0$ can be written as $\varphi(a_1,a_2)=\tilde{\varphi}(a_1a_2)$
for some $K$-linear $\tilde{\varphi}\colon A^2\to B$.
\end{definition}

The notion is from Bre\v{s}ar, Gra\v{s}i\'c and S\'anchez Ortega, who
introduced it for matrix algebras
\cite[pp.~1486--1498]{BresarGrasicSanchezOrtega2009} and state it in the form
used here in the later paper of Brice and Huang
\cite[pp.~326--327]{BriceHuang2015}.
It was built for zero product preserving
linear maps on Banach algebras and on matrix algebras under the standard,
Lie and Jordan products \cite[p.~327]{BriceHuang2015}.
Brice and Huang give a tensor criterion for the class of
Definition~\ref{def:local:zpd}.

\begin{theorem}[{\cite[Theorem~2.3, p.~328]{BriceHuang2015}}]\label{thm:local:zpdclass}
Let $(A,\mu)$ be a $K$-algebra, with $\mu\colon A\otimes A\to A$ its
multiplication map and $T_\mu$ the span of all $a_1\otimes a_2$ with
$a_1a_2=0$. Then $(A,\mu)$ is zero product determined if and only if
$T_\mu=\ker\mu$.
\end{theorem}

The criterion of Theorem~\ref{thm:local:zpdclass} runs the class well. A
direct sum of $K$-algebras belongs to the class exactly when each summand
does \cite[Theorem~3.1, p.~331]{BriceHuang2015}. An algebra tensor product
of two zero product determined algebras stays in the class when $K$ is a
field or both multiplications are surjective
\cite[Theorem~4.3, pp.~331--332]{BriceHuang2015}.

An element can carry the notion instead of the map. Zhu calls $Z$ an
all-derivable point of an operator algebra when every continuous derivable map at
$Z$ is a derivation, and identifies such points
\cite[Abstract]{Zhu2007AllDerivable}.

Benkovic and Grasic settle when the zero-product condition determines the map.
They show a unital algebra with a nontrivial idempotent is zero product determined
exactly when both Peirce subalgebras are, and then characterize the generalized
derivations determined by their action on zero products
\cite[pp.~347--368]{BenkovicGrasic2014Generalized}.

Further work in the operator and zero-product line, held in the collection:
\noindent\emph{Generalized skew derivations with nilpotent values on left ideals of rings and Banach algebras}~\cite{LibGeneralizedSkewDerivationsWith2014}.
\noindent\emph{Aggregate special functions to approximate permuting tri-homomorphisms and permuting tri-derivations associated with a tri-additive $\psi$ -functional inequality in Banach algebras}~\cite{LibAggregateSpecialFunctionsTo2024}.
\noindent\emph{Additive maps on some operator algebras behaving like $\alpha$ $\beta$ -derivations or generalized $\alpha$ $\beta$ -derivations at zero-product elements}~\cite{LibAdditiveMapsOnSome2014}.
\noindent\emph{Ring isomorphisms and linear or additive maps preserving zero products on nest algebras}~\cite{LibRingIsomorphismsAndLinear2004}.

The Jordan version has the same shape. An and Hou call $Z$ a Jordan
full-derivable point when every additive map Jordan derivable at $Z$ is a Jordan
derivation, and prove that the unit is one on a prime ring under mild conditions,
on a triangular ring, on a factor von Neumann algebra and on a nest algebra
\cite[Abstract, p.~753]{AnHou2010JordanFullDerivable}.

An additive map of a ring that behaves like a derivation at every idempotent is shown to be a derivation, which is the idempotent form of the zero-product characterization of Chapter 8~\cite{WlAdditiveMapsOnRingsBehavingLik2011}.

An additive map derivable, or Jordan derivable, at the zero points of a ring is shown to be a derivation under the stated conditions~\cite{WlAdditiveMapsDerivableOrJordanD2010}.

For a subalgebra of $B(H)$ with the identity, a linear map into $B(H)$ that is a generalized derivable mapping at the zero point is characterised on the reflexive operator algebras~\cite{WlGeneralizedDerivableMappingsAt2019}.

A map of a $C^*$-algebra into its dual space that is derivable at a single point is shown to be a derivation, which extends derivability from one point to the whole algebra~\cite{WlDualSpaceValuedMappingsOnCAlge2019}.

For a standard operator algebra on a complex Banach space of dimension at least two, a linear map that is anti-derivable at zero is shown to be an anti-derivation plus a term the chapter's zero-product method identifies~\cite{WlAntiDerivableLinearMapsAtZeroO2019}.

Moln\'ar studies the derivations of the standard operator algebras on a Banach space that are locally inner, that is, inner at every point~\cite{WlLocallyInnerDerivationsOfStand1996}.

Qi and Hou show that for a $J$-subspace lattice $\mathcal{L}$ on a Banach space with $\dim K\neq 2$ for every $K$ in $J(\mathcal{L})$, every linear map on the algebra of finite-rank operators of $\mathrm{Alg}\,\mathcal{L}$ that is Lie derivable at zero has the standard form $A\mapsto\delta(A)+\varphi(A)$ with $\delta$ a generalized derivation and $\varphi$ a centre-valued linear map, and characterise the Lie derivable maps at zero on $\mathrm{Alg}\,\mathcal{L}$ as well~\cite{WlLinearMapsLieDerivableAtZeroOn2010}.

Zivari-Kazempour and Ghahramani characterise the continuous linear maps $\delta$ of a unital Banach algebra into a Banach bimodule that are derivable, or generalized derivable, at the pairs $a,b$ with $ab=ba=w$, where $w$ is a separating point in the centre or an idempotent~\cite{WlDerivableMapsAtCommutativeProd2024}.

For a von Neumann algebra $M$ with no central summands of type $I_1$ or $I_2$, Qi, Ji and Hou show that an additive map $L$ satisfies the $\xi$-Lie derivation rule at zero, $L([A,B])=[L(A),B]+[A,L(B)]$ whenever $[A,B]=0$, exactly when it has the stated form involving an element of the centre, an additive map into the centre and an additive derivation~\cite{WlCharacterizationOfAdditiveMaps2015}.
